\section{The arithmetic quotient}
\label{sec:arithmetic-quotient}

We first identify the quotient through which the multiplier and the Kummer
cocycle of~$2$ factor. Put
\[
 G=\widehat{\mathrm{GT}},\quad H=G_{\mathbb Q},\quad
 N=\Lambda _0,\quad
 K_2=\mathbb Q(\mu _\infty,2^{1/\infty}),\quad K=G_{K_2},
 \qquad U=\widehat{\mathbb Z}^{\times}.
\]
The embedding of $H$ in $G$ is understood throughout. We use the decomposition
$G=NH$ of \cite[Theorem~2]{H17}, the cocycle identity
\[
 \rho _2(gg')=\rho _2(g)+\lambda(g)\rho _2(g'),
\]
and the fact that $\rho _2$ restricts to the Kummer cocycle of the positive
coherent roots of~$2$; see \cite[Definition~5.1]{NS2000} and
\cite[Lemma~7]{H17}. Define the additive quadratic character
$\epsilon _8:U\to\mathbb F _2$ by
\[
 \epsilon _8(a)=
 \begin{cases}
 0,&a\equiv 1,7\pmod 8,\\
 1,&a\equiv 3,5\pmod 8.
 \end{cases}
\]

\begin{proposition}\label{prop:arithmetic-affine}
The subgroup $N$ is closed and normal in $G$, and $N\cap H=K$.
The map $(\lambda,\rho _2)$ identifies $G/N$ with
\[
 A=\{(a,r)\in U\ltimes\widehat{\mathbb Z}:
                 r\bmod 2=\epsilon _8(a)\},
 \qquad
 (a,r)(b,s)=(ab,r+as).
\]
This identifies $G/N$ canonically with $\operatorname{Gal}(K_2/\mathbb Q)$.
In particular $\rho _2\bmod 2=\epsilon _8\circ\lambda$ on $G$.
\end{proposition}

\begin{proof}
The cocycle identity makes $(\lambda,\rho _2)$ a continuous homomorphism to
the affine semidirect product. Its kernel is $N$, so $N$ is closed and normal.
On $H$ the simultaneous kernel fixes every root of unity and every coherent
root of~$2$, and is consequently $K$.

We need the precise Kummer image over $\mathbb Q^{\mathrm{ab}}
=\mathbb Q(\mu _\infty)$. It is a closed subgroup $J$ of
$\widehat{\mathbb Z}$. For a positive integer $d$, the inclusion
$J\subset d\widehat{\mathbb Z}$ is equivalent to
$2^{1/d}\in\mathbb Q^{\mathrm{ab}}$. This happens exactly for $d=1,2$.
Indeed $X^d-2$ is Eisenstein; if $d>2$, the field generated by its positive
real root is real and misses a nonreal conjugate, so is not normal over
$\mathbb Q$. Every finite subfield of an abelian Galois extension is normal.
Closed subgroups of $\widehat{\mathbb Z}$ have the form
$\prod_p p^{e_p}\mathbb Z_p$ with $e_p\in\mathbb N\cup\{\infty\}$,
and the inclusions in $d\widehat{\mathbb Z}$ determine all the $e_p$.
Thus $J=2\widehat{\mathbb Z}$.

Given $g\in G$, choose $h\in H$ with $\lambda(h)=\lambda(g)$.
The element $gh^{-1}$ has multiplier one, and
\cite[Lemma~7(c)]{H17} gives $\rho _2(gh^{-1})\in2\widehat{\mathbb Z}$.
The cocycle identity therefore gives
$\rho _2(g)\equiv\rho _2(h)\pmod2$. On $H$ this parity is the action
on $\sqrt2=\zeta _8+\zeta _8^{-1}$, namely $\epsilon _8(\lambda(h))$.
For each fixed multiplier, the Kummer values on $H$ fill the corresponding
coset of $2\widehat{\mathbb Z}$, by the preceding calculation. Thus the
image of $H$ is exactly $A$. The decomposition $G=NH$ shows that $G$ has
the same image. Finally the continuous bijection $H/K\to G/N$ is a
homeomorphism, since its source is compact and its target Hausdorff.
\end{proof}

\begin{proposition}\label{prop:arithmetic-commutators}
Let $T=\{(1,2t):t\in\widehat{\mathbb Z}\}\subset A$.
Then
\[
 \overline{[A,A]}=T,\qquad
 \overline{[A,T]}=\{(1,4t):t\in\widehat{\mathbb Z}\}.
\]
Every finite $p$-group quotient of $A$ is abelian when $p$ is odd.
Every finite quotient of $A$ which is a $2$-group of nilpotency class at
most two has commutator subgroup of order at most two.
\end{proposition}

\begin{proof}
The quotient $A/T=U$ is abelian. Choose $a\in U$ whose $2$-adic
component is~$3$ and whose odd-primary components are~$1$.
Then $(a,1)\in A$, and its commutator with $c=(-1,0)$ is $(1,2)$.
The closed subgroup generated by this translation is $T$, proving the first
equality. Conjugation by $(a,r)$ multiplies a translation coordinate by~$a$.
Thus $[A,T]$ consists of coordinates generated by $(a-1)2t$.
They are all divisible by~$4$ at the prime~$2$, and $a=-1$ produces
$(1,-4)$, whose closed cyclic subgroup is $4\widehat{\mathbb Z}$.
This proves the second equality.

The element $c$ has order two and acts on $T$ by inversion. It maps to the
identity in an odd-primary finite group, so the image of $T$ in such a
group has exponent dividing two and is trivial. The quotient is therefore
abelian. A quotient of nilpotency class at most two kills $[A,T]$, and its
commutator subgroup is consequently a quotient of
$T/\overline{[A,T]}\simeq C_2$.
\end{proof}

\section{Characters and the cyclotomic abelianization}
\label{sec:characters-abelianization}

The decomposition $G=NH$ separates the characters obtained from the multiplier
from the characters that could distinguish $G$ from $H$. All characters
in this section are continuous, and all abelianizations are profinite.

\begin{lemma}\label{lem:character-gluing}
Let $G=NH$ be a profinite group with $N$ closed and normal and $H$ closed,
and put $K=N\cap H$. Let $M$ be a finite abelian group with trivial action.
An invariant character $\psi:N\to M$ and a character $\eta:H\to M$
extend to a character of $G$ if and only if $\psi|_K=\eta|_K$.
The extension is unique and is given by
\[
 \chi(nh)=\psi(n)+\eta(h).
\]
\end{lemma}

\begin{proof}
Necessity follows by restriction. If $nh=n'h'$, then
$n'^{-1}n=h'h^{-1}\in K$, so the proposed two values coincide.
The identity
\[
 (n_1h_1)(n_2h_2)=n_1(h_1n_2h_1^{-1})h_1h_2
\]
and invariance of $\psi$ prove additivity. The multiplication map
$N\times H\to G$ is a continuous surjection from a compact space to a
Hausdorff space and hence is a quotient map. The displayed function on
$N\times H$ is continuous and constant on its fibres, proving continuity
on $G$. Uniqueness follows from $G=NH$.
\end{proof}

In our arithmetic situation define
\[
 C_N=\overline{\langle[g,n],k:g\in G,\ n\in N,\ k\in K\rangle},
 \qquad B=N/C_N.
\]
The group $B$ is abelian, since $C_N$ contains $[N,N]$; it is also the
largest quotient of $N$ on which conjugation by $G$ and the image of $K$
are both trivial.

\begin{theorem}\label{thm:abelian-splitting}
There is a canonical isomorphism of profinite abelian groups
\[
 G^{\mathrm{ab}}\simeq U\times B
\]
under which the multiplier is the first projection. Moreover
\[
 C_N=N\cap\overline{[G,G]},\qquad
 B\simeq\bigl(G/\overline{\langle\!\langle H\rangle\!\rangle}\bigr)^{\mathrm{ab}}.
\]
For every prime $p$, put
\[
 X_p=H^1(N,\mathbb F_p)^G,
 \qquad r_p:X_p\longrightarrow H^1(K,\mathbb F_p).
\]
Restriction gives canonical identifications
\begin{align*}
 \ker r_p&\simeq\operatorname{Hom}_{\mathrm{cts}}(B,\mathbb F_p),\\
 H^1(G,\mathbb F_p)&\simeq
 \lambda^*H^1(U,\mathbb F_p)\oplus\ker r_p.
\end{align*}
Consequently $\ker r_p=0$ if and only if every continuous
$\mathbb F_p$-character of $G$ factors through the multiplier.
The multiplier identifies $G^{\mathrm{ab}}$ with $U$ if and only if
$N\subset\overline{[G,G]}$, and this holds if and only if $\ker r_p=0$
for every prime~$p$.
\end{theorem}

\begin{proof}
Kronecker--Weber identifies $H^{\mathrm{ab}}$ with $U$ through the
cyclotomic character. Since $K$ fixes $\mathbb Q^{\mathrm{ab}}$, it lies
in $\overline{[H,H]}$ and hence in $\overline{[G,G]}$.
The map $N\to G^{\mathrm{ab}}$ therefore factors through $B$.
Conversely the rule
\[
 nh\longmapsto(\lambda(h),n\bmod C_N)
\]
is well-defined: changing the decomposition changes $n$ by an element of
$K$, which is zero in $B$. It is continuous by the compact quotient
argument of Lemma~\ref{lem:character-gluing}. It is a homomorphism because
conjugation by $H$ is trivial on $B$. The maps
$H^{\mathrm{ab}}=U\to G^{\mathrm{ab}}$ and $B\to G^{\mathrm{ab}}$
give the inverse of the resulting homomorphism
$G^{\mathrm{ab}}\to U\times B$: both composites are the identity on the
images of $N$ and $H$, which generate the relevant groups.
This proves the splitting and the formula for $C_N$.
Killing $H$ and then abelianizing imposes exactly the relations defining
$C_N$, which proves the second description of~$B$.

Apply Lemma~\ref{lem:character-gluing} with $\eta=0$ to identify the
restriction kernel. Every character of $H$ is the pullback of a unique
character of $U$ and vanishes on $K$. Subtracting that pullback from a
character of $G$ gives the asserted direct sum, or it follows directly
from the abelian splitting. Finally a nonzero profinite abelian group has
a nonzero finite abelian quotient, and that quotient has a quotient
$C_p$ for some prime $p$. Thus $B=0$ if and only if it has no character
to $\mathbb F_p$ for any prime~$p$.
\end{proof}

The expression ``the abelianization is $U$'' in this theorem specifies
the multiplier map. The kernel assertion does not follow from an
unmarked abstract isomorphism of profinite abelian groups.

\subsection{The transgression describes the image}

The kernel of restriction and the image of restriction involve different
obstructions. The next proposition states both in one exact sequence.
Write
\[
 Y_p=H^1(K,\mathbb F_p)^H.
\]
Restriction $r_p$ lands in $Y_p$, and $A=G/N=H/K$ acts on $K$ through
conjugation by $H$.

\begin{proposition}\label{prop:transgression-comparison}
The transgression
$d_H:Y_p\to H^2(A,\mathbb F_p)$ is injective and identifies $Y_p$ with
\[
 \ker\bigl(H^2(A,\mathbb F_p)\longrightarrow H^2(H,\mathbb F_p)\bigr).
\]
There is an exact sequence
\[
 0\longrightarrow\operatorname{Hom}_{\mathrm{cts}}(B,\mathbb F_p)
 \longrightarrow X_p
 \xrightarrow{\ d_Hr_p\ }
 \ker\bigl(H^2(A,\mathbb F_p)\longrightarrow H^2(G,\mathbb F_p)\bigr)
 \longrightarrow0.
\]
Thus, under the injection $d_H$, the image of restriction is precisely
the inflation kernel for~$G$.
\end{proposition}

\begin{proof}
For completeness, we describe the central-extension form of the
five-term argument. Given a conjugation-invariant character
$\psi:N\to\mathbb F_p$, its pushout extension is
\[
 E_\psi=(\mathbb F_p\times G)/
       \{(-\psi(n),n):n\in N\}.
\]
The subgroup being divided out is closed and normal, by the invariance
of $\psi$. The quotient is a central extension of $A$ by $\mathbb F_p$;
its class is $d_G\psi$, with a consistent choice of the transgression
sign. A splitting of this extension is equivalent to a character of
$G$ extending $\psi$: in one direction subtract the splitting from the
natural lift $g\mapsto[0,g]$, and in the other use the character to
define a splitting independent of the lift from~$A$.

More generally, a central extension of $A$ becomes trivial after
inflation to $G$ if and only if it admits a continuous homomorphism
$G\to E$ lifting $G\to A$. Restricting such a lift to $N$ gives an
invariant character $\psi$, and the map
$[z,g]\mapsto z\,\widetilde g$ identifies its pushout with $E$.
This proves exactness at $X_p$ and identifies the image of $d_G$ with
the indicated inflation kernel. The construction respects sums of
characters and Baer sums of central extensions, so these are linear
maps.

Apply the same argument to $H$ and $K$. A character of $K$ in the
kernel of $d_H$ extends to $H$. Every character of $H$ factors through
$U$ and therefore vanishes on $K$, so $d_H$ is injective. The preceding
lifting argument gives its image. Finally pulling the pushout extension
back from $G/N$ along $H/K$ is the same construction with $\psi|_K$.
Thus $d_G=d_Hr_p$. Theorem~\ref{thm:abelian-splitting} identifies the
kernel of $d_G$, proving the displayed exact sequence.
\end{proof}

\section{The arithmetic invariant characters}
\label{sec:arithmetic-target}

The target allowed by Galois invariance is already infinite dimensional.
Its size does not determine which characters extend across $K\subset N$.

\begin{theorem}\label{thm:arithmetic-target-dimension}
For every prime $p$,
\[
 \dim_{\mathbb F_p}H^1(G_{K_2},\mathbb F_p)^{G_{\mathbb Q}}=\aleph_0.
\]
For $p=2$, an explicit independent family is given, for the odd rational
primes $q$, by the quadratic extensions $K_2(q^{1/4})/K_2$.
Equivalently these are the Kummer classes of $\sqrt q$ in
$K_2^\times/(K_2^\times)^2$.
\end{theorem}

\begin{proof}
For an odd prime $q$, the square of the positive fourth root $q^{1/4}$
belongs to $\mathbb Q^{\mathrm{ab}}\subset K_2$. The extension in the
statement is therefore at most quadratic. It is normal over $\mathbb Q$,
because its other conjugates are $i^j q^{1/4}$ and $i\in K_2$.
Its character, if nonzero, is invariant under $H$, because conjugation
acts trivially on a group of order two.

Suppose that a nonempty finite product of these characters is zero.
Let $a$ be the product of the corresponding distinct odd primes. Then
$a^{1/4}\in K_2$. The extension
$\mathbb Q^{\mathrm{ab}}(a^{1/4})/\mathbb Q^{\mathrm{ab}}$ is a
nontrivial quadratic extension. To see nontriviality, $X^4-a$ is
Eisenstein at any prime dividing $a$, and its positive real root
generates a nonnormal real field, which cannot lie in
$\mathbb Q^{\mathrm{ab}}$. By Proposition~\ref{prop:arithmetic-affine},
$\operatorname{Gal}(K_2/\mathbb Q^{\mathrm{ab}})=2\widehat{\mathbb Z}$
has exactly one subgroup of index two. Its quadratic subfield is
$\mathbb Q^{\mathrm{ab}}(2^{1/4})$. The two quadratic fields are
therefore equal, so the product $a^{1/4}2^{1/4}=(2a)^{1/4}$ lies in
$\mathbb Q^{\mathrm{ab}}$: their common nontrivial automorphism negates
both factors. But $X^4-2a$ is Eisenstein at~$2$, and its positive root
again generates a real nonnormal field. This contradiction proves
nontriviality and independence of the whole quadratic family.

Now let $p$ be odd and let $m\geq1$. Put
$P=\operatorname{UT}_3(\mathbb F_p)^m$, a direct product of upper
unitriangular matrix groups. Its centre and commutator subgroup agree
and are isomorphic to $(C_p)^m$, and its abelianization is
$(C_p)^{2m}$. By \v{S}afarevi\v{c}'s theorem, in the precise form
\cite[Theorem~1, p.~1]{SchmidtWingberg}, there is a finite Galois
extension $E/\mathbb Q$ with group~$P$.
Let $L=E\cap K_2$. Its Galois group over $\mathbb Q$ is a finite
$p$-group quotient of $A$, and is abelian by
Proposition~\ref{prop:arithmetic-commutators}. Hence
$L\subset E^{[P,P]}$. The reverse inclusion follows from
$E^{[P,P]}\subset\mathbb Q^{\mathrm{ab}}\subset K_2$. Thus
\[
 \operatorname{Gal}(EK_2/K_2)=\operatorname{Gal}(E/L)=[P,P]
                  \simeq(C_p)^m.
\]
This subgroup is central in $\operatorname{Gal}(EK_2/\mathbb Q)$,
because $[P,P]$ is central in~$P$. Its characters give $m$ independent
members of $Y_p$.

Finally the algebraic field $\overline{\mathbb Q}$ is countable.
Every cyclic degree-$p$ extension of $K_2$ inside it is generated by
one element, so there are at most countably many such extensions.
Each has finitely many characters. This gives the countable upper
bound, while the preceding arguments give arbitrarily large finite
independent families for every~$p$.
\end{proof}

\section{Infinitely many invariant quadratic characters of the normalized group}
\label{sec:quadratic-dimension}

Ihara's Kummer cocycles give explicit invariant quadratic characters and
determine the dimension of their restriction image. They do not settle
the kernel of restriction. We write
$\rho_d^I$ for these cocycles to distinguish their construction from the
frame cocycle used to define~$N$.

Let $d$ be a positive integer and put
\[
 F=\widehat F_2=\langle x,y\rangle,\qquad
 H_d=\ker(F\longrightarrow\mathbb Z/d\mathbb Z),
 \quad x\longmapsto1,\quad y\longmapsto0.
\]
The Schreier basis of $H_d$ consists of $x^d$ and
$y_j=x^jyx^{-j}$ for $0\leq j<d$. Let
$\ell_d:H_d\to\widehat{\mathbb Z}$ send $y_0$ to~$1$ and the other
basis elements to~$0$. For $g=(\lambda,f)\in\GTz$ define
\[
 \rho_d^I(g)=-\ell_d(f).
\]
This is Ihara's cocycle $-\Psi_d^{(0)}$; the construction and its
identification with the Kummer cocycle of $d$ on $G_{\mathbb Q}$ are
recalled in \cite[\S2.2]{Furusho2003}, following Ihara's theory
\cite{Ihara1999}. We use that arithmetic identification as a published
input. The following direct verification fixes the cocycle convention.

\begin{lemma}\label{lem:ihara-cocycle}
For $g,g'\in\GTz$,
\[
 \rho_d^I(gg')=\rho_d^I(g)+\lambda(g)\rho_d^I(g').
\]
If $n\in N$, then
\[
 \rho_d^I(gng^{-1})=\lambda(g)\rho_d^I(n).
\]
\end{lemma}

\begin{proof}
The automorphism $\Phi_g:x\mapsto x^\lambda$,
$y\mapsto f^{-1}y^\lambda f$ preserves $H_d$. On its abelianization,
\[
 [x^d]\longmapsto\lambda[x^d],\qquad
 [y_j]\longmapsto\lambda[y_{\lambda j\bmod d}].
\]
The two conjugating contributions of $f$ cancel. Replacing an exponent
$\lambda j$ by its residue modulo $d$ only conjugates by a power of
$x^d$, which is also invisible in the abelianization. Since $\lambda$
is a unit modulo $d$, index zero has no preimage other than zero.
Consequently $\ell_d\Phi_g=\lambda\ell_d$.
With the convention $\Phi_{gg'}=\Phi_g\Phi_{g'}$, the second coordinate
is $f_{gg'}=f_g\Phi_g(f_{g'})$. Applying $-\ell_d$ proves the first
formula. Applying it twice to $gng^{-1}$, and using $\lambda(n)=1$,
cancels the terms involving $\rho_d^I(g)$ and proves the second.
\end{proof}

\begin{lemma}\label{lem:cyclic-cover-symmetry}
Let $F=\widehat F_2=\langle x,y\rangle$. Suppose a continuous
automorphism $\phi$ satisfies
\[
 \phi(x)=x,\qquad \phi(y)=f^{-1}yf,\qquad
 f\in\overline{[F,F]},
\]
and $\phi(xy)$ is conjugate to $xy$. For a positive integer $n$, use
the free basis
\[
 X=x^n,\qquad Y_j=x^jyx^{-j}\quad(0\leq j<n)
\]
of $H_n$. If $[f]=a[X]+\sum_{j=0}^{n-1}c_j[Y_j]$ in
$H_n^{\mathrm{ab}}$, then
\[
 a=0,\qquad \sum_{j=0}^{n-1}c_j=0,\qquad
 c_j=c_{-j\bmod n}\quad(0\leq j<n).
\]
\end{lemma}

\begin{proof}
The map $H_n^{\mathrm{ab}}\to F^{\mathrm{ab}}$ takes $[X]$ to
$n[x]$ and every $[Y_j]$ to $[y]$. Hence $na=0$ and $\sum c_j=0$.
Multiplication by $n$ is injective on $\widehat{\mathbb Z}$, so $a=0$.
The automorphism preserves $H_n$, fixes $X$, and sends
\[
 Y_i\longmapsto f_i^{-1}Y_if_i,\qquad f_i=x^ifx^{-i}.
\]
Conjugation by $x$ cyclically permutes the $[Y_j]$ in the abelianization
and fixes $[X]$, so $[f_i]=\sum_jc_{j-i}[Y_j]$, with indices modulo $n$.

The boundary word
\[
 P=Y_0Y_1\cdots Y_{n-1}X=(yx)^n
\]
is a power of a conjugate of $xy$, and therefore
$\phi(P)=tPt^{-1}$ for some $t\in F$. Multiplying $t$ on the right
by a suitable integer power of $yx$ makes $t$ belong to $H_n$ without
changing this equality: $yx$ centralizes $P$ and generates $F/H_n$.

Work in the quotient of $H_n$ of nilpotency class two, and use the
commutator convention $[u,v]=uvu^{-1}v^{-1}$. The commutator pairing
identifies its degree-two group with the exterior square of
$H_n^{\mathrm{ab}}$ over $\widehat{\mathbb Z}$. In particular the
commutators of distinct basis elements are independent. This independence
can also be checked in finite unitriangular $3$-by-$3$ matrix groups:
map one chosen pair of generators to the two elementary generators and
all other generators to the identity, for every finite modulus.
Using the same symbols for degree-one vectors, the boundary identity gives
\[
 \sum_{i=0}^{n-1}[Y_i,[f_i]]
       =\left[[t],X+\sum_{i=0}^{n-1}Y_i\right].
\]
Here each outer bracket is the bilinear commutator pairing; $[f_i]$
and $[t]$ denote degree-one vectors. Write
$[t]=t_XX+\sum_i t_iY_i$. The coefficient of $[X,Y_i]$ on the left
is zero because every $[f_i]$ has zero $X$-coordinate. On the right it
is $t_X-t_i$. Thus every $t_i=t_X$, and the right side is zero.
For $i<j$, the left side has coefficient $c_{j-i}-c_{i-j}$ at
$[Y_i,Y_j]$. Taking $i=0$ gives the asserted symmetry.
\end{proof}

\begin{theorem}\label{thm:all-kummer-parity}
For every positive integer $d$, Ihara's integral Kummer cocycle is even
on the multiplier-one subgroup of the gentle group:
\[
 \rho_d^I\bigl(\ker(\lambda:\widehat{\mathrm{GT}}_0\to U)\bigr)
       \subset2\widehat{\mathbb Z}.
\]
On the whole gentle group it satisfies
\[
 \rho_d^I\bmod2=\epsilon_d\circ\lambda,
\]
where $\epsilon_d:U\to\mathbb F_2$ is the additive quadratic
Dirichlet character of $\mathbb Q(\sqrt d)$.
\end{theorem}

\begin{proof}
The two cycle equations at multiplier one give the positive peripheral
preservation in Lemma~\ref{lem:cyclic-cover-symmetry}. Apply that lemma
with $n=2d$. Pairing indices $j$ and $-j$ gives
\[
 0=\sum_{j=0}^{2d-1}c_j
   =c_0+c_d+2\sum_{j=1}^{d-1}c_j.
\]
Under $H_{2d}\subset H_d$, only $Y_0$ and $Y_d$ among these basis
elements have nonzero $y$-coordinate in $H_d^{\mathrm{ab}}$, and that
coordinate is one for each. Indeed $Y_{j+d}=x^dY_jx^{-d}$.
The element $x^{2d}$ has zero $y$-coordinate as well. Consequently
\[
 \rho_d^I=-\ell_d(f)=-(c_0+c_d)\in2\widehat{\mathbb Z}.
\]
For a general gentle element $g$, choose $h\in G_{\mathbb Q}$ with
$\lambda(h)=\lambda(g)$. The cocycle identity and the first assertion
give $\rho_d^I(g)\equiv\rho_d^I(h)\pmod2$. The latter is the action
of $h$ on $\sqrt d$, and hence is $\epsilon_d(\lambda(g))$.
\end{proof}

\begin{theorem}\label{thm:quadratic-dimension}
The vector space of invariant continuous quadratic characters of the
normalized subgroup has dimension
\[
 \dim_{\mathbb F_2}H^1(\Lambda_0,\mathbb F_2)^{\widehat{\mathrm{GT}}}
 =\aleph_0.
\]
For every positive integer $d$ the formula
\[
 \psi_d^N=\frac{\rho_d^I|_N}{2}\bmod2
\]
defines a continuous invariant quadratic character. For the odd rational
primes $q$, the characters $\psi_q^N$ are linearly independent and
restrict to the characters of $K_2(q^{1/4})/K_2$. Thus
\[
 \dim_{\mathbb F_2}\operatorname{im}r_2=\aleph_0.
\]
More precisely, for distinct odd primes $q_1,\ldots,q_m$, let
$R:N\to\mathbb Z_2^m$ be the vector of the $2$-adic components of
$\rho_{q_1}^I,\ldots,\rho_{q_m}^I$. Then
\[
 R(K)=R(N)=2\mathbb Z_2^m,
\]
and restriction is injective on the $m$-dimensional character space
obtained by pullback from this lattice.
\end{theorem}

\begin{proof}
Theorem~\ref{thm:all-kummer-parity} permits division by two on $N$.
Since $\lambda=1$ there, the result is a continuous homomorphism.
The conjugation formula of Lemma~\ref{lem:ihara-cocycle} and the fact
that every multiplier is odd make its reduction modulo two invariant.
On $K$, the character $\psi_q^N$ is the ordinary fourth-root Kummer
character because $K$ fixes $\sqrt q$. Independence follows from
Theorem~\ref{thm:arithmetic-target-dimension}.

The vector $R$ is a continuous homomorphism. On $K$ it is the vector of ordinary Kummer
characters of the $q_j$. All their values are even, since $K$ fixes
$\sqrt{q_j}$. Modulo four, half their values are precisely the
independent quadratic characters of
Theorem~\ref{thm:arithmetic-target-dimension}. Thus the image of $R(K)$
modulo $4\mathbb Z_2^m$ is all of
$2\mathbb Z_2^m/4\mathbb Z_2^m$.

A closed additive subgroup of $\mathbb Z_2^m$ is a
$\mathbb Z_2$-submodule, because integer multiplication extends by
continuity to $2$-adic scalar multiplication. A closed submodule
$M\subset2\mathbb Z_2^m$ whose reduction modulo four is onto equals
$2\mathbb Z_2^m$. Explicitly, approximate any vector modulo four by
an element of $M$, then correct the error modulo eight by twice an
element of $M$, and continue. The convergent sum belongs to $M$ by
closedness. Apply this to the compact image $R(K)$.

The parity theorem gives $R(N)\subset2\mathbb Z_2^m$, and the reverse
inclusion follows from $K\subset N$. Hence $R(N)=R(K)$.
Conjugation acts on this lattice by the odd scalar $\lambda(g)$, so
all its quadratic characters are invariant. Its character space has
dimension $m$, and restriction to $K$ is injective because $R(K)=R(N)$.

There is also a countable upper bound. The finitely generated free
profinite group $\widehat F_2$ has a countable base, as there are only
countably many finite quotients with specified generator images.
Thus $U\times\widehat F_2$, its closed subspace $G$, and the closed
subgroup $N$ are compact metrizable spaces. Such a space has countably
many clopen subsets: each is a finite union of members of a countable
base. A continuous function to $\mathbb F_2$ is determined by its
zero fibre. Consequently $H^1(N,\mathbb F_2)$, and hence its invariant
subspace, is countable. The lower bounds for arbitrary $m$ prove the
dimension formula.

The same upper bound applies to the restriction image, while the
independent quartic restrictions give its countably infinite dimension.
\end{proof}

The dimension computation identifies the explicit quartic subspace. To determine its complete extent, and to normalize all rational Kummer coordinates simultaneously, we next use the distribution theorem.
